% Source: https://artofproblemsolving.com/community/c3h1077508_how_to_quickly_find_sqaure_roots
% Stable identity: aops-topic-1077508-q01
\textbf{(a)}
\[
x'-\sqrt N=\dfrac{x^2+N-2x\sqrt N}{2x}=\dfrac{(x-\sqrt N)^2}{2x}\ge0 .
\]

\textbf{(b)} $x_1=\dfrac{1}{2}\left(3+\dfrac{10}{3}\right)=\dfrac{19}{6}$ and
\[
x_2=\dfrac{1}{2}\left(\dfrac{19}{6}+\dfrac{60}{19}\right)=\dfrac{361+360}{228}=\dfrac{721}{228}\approx3.1622807 .
\]
By (a), $x_1-\sqrt{10}=\dfrac{(3-\sqrt{10})^2}{6}$. Since $3<\sqrt{10}<3.17$, we get $0<x_1-\sqrt{10}<\dfrac{0.17^2}{6}<0.005$. Applying (a) again, with $x_1>3$,
\[
0<x_2-\sqrt{10}=\dfrac{(x_1-\sqrt{10})^2}{2x_1}<\dfrac{0.005^2}{6}<5\cdot10^{-6}<10^{-5}.
\]
(Positivity: $x_1\ne\sqrt{10}$ because $x_1$ is rational.) So two steps of the method already give $\sqrt{10}$ to five decimal places: $\sqrt{10}\approx3.16228$. Each step roughly doubles the number of correct digits.
